% !PRINTED 469
\weberpage{469}

\[
\sigma_a = (\varrho, \varrho_a), \qquad a = 0, 1, 2 \ldots \mu - 1,
\]
wobei unter $\sigma_0$ oder $\sigma$ die sogenannte \gesperrt{identische Substitution} $(\varrho, \varrho)$ verstanden wird, die darin besteht, dass $\varrho$ durch sich selbst ersetzt wird, also alle Zahlen in $\Omega(\varrho)$ ungeändert bleiben.

Wenn wir in einer beliebigen der Wurzeln (3)
\[
\varrho_h = \Theta_h(\varrho)
\]
die Substitution $\sigma_a$ ausführen, so geht $\varrho_h$ in eine andere Wurzel $\varrho_k$ über, die bestimmt ist durch
\[
\varrho_k = \Theta_h(\varrho_a) = \Theta_h \Theta_a(\varrho) = \Theta_k(\varrho).
\]
Es ist also $\varrho_k$ dieselbe Function von $\varrho_a$, wie $\varrho_h$ von $\varrho$.

Eine beliebige Grösse $\omega = \varphi(\varrho)$ des Körpers $\Omega(\varrho)$ geht durch die Substitution $\sigma_a$ in $\omega_a = \varphi(\varrho_a)$ über. Drücken wir $\omega$ durch $\varrho_h$ und $\omega_a$ durch $\varrho_k$ aus:
\[
\omega = \psi(\varrho_h), \qquad \omega_a = \psi(\varrho_k),
\]
so sieht man, dass $\omega$ in $\omega_a$ übergeht durch die Substitution
\weq{9}{\sigma_a = (\varrho_h, \varrho_k) = (\varrho_h, \Theta_h \Theta_a(\varrho)),}
und dies ist nur ein anderer Ausdruck für die Substitution $(\varrho, \varrho_a)$.

Hierin ist $\varrho_h$ eine beliebige Wurzel von $g(t)$ und das zugehörige $\varrho_k$ ist durch $\sigma_a$ bestimmt. Man kann aber bei gegebenem $\sigma_a$ auch $\varrho_k$ beliebig annehmen, und das zugehörige $\varrho_h$ bestimmen, indem man in $\Theta_h(\varrho_a)$ den Index $h$ die Werthe $0, 1, 2 \ldots \mu - 1$ durchlaufen lässt, wobei jede Wurzel $\varrho_k$ einmal zum Vorschein kommt. Also:

\gesperrt{3. Jede Substitution $\sigma_a$ kann in der Form $(\varrho_h, \varrho_k)$ dargestellt werden, worin entweder $\varrho_h$ oder $\varrho_k$ eine beliebig gegebene der $\mu$ Wurzeln $\varrho$ ist, während die andere durch $\sigma_a$ völlig bestimmt ist.}

% [sic] Der Druck bricht den Satz grammatisch nach „in die gleiche“ ab; gemeint ist offenbar „in die gleiche Grösse übergehen“. Der Wortlaut bleibt quellentreu erhalten.
Die Anzahl aller von einander verschiedenen Substitutionen $\sigma_a$ ist also, die identische Substitution mit gerechnet, gleich dem Grade des Körpers $\Omega(\varrho)$, den wir oben schon mit $\mu$ bezeichnet haben. Jede dieser Substitutionen führt die Gesammtheit der Grössen des Körpers $\Omega(\varrho)$ in sich selbst über, so dass jede in eine bestimmte andere übergeht, und niemals zwei verschiedene Grössen in die gleiche [Formel (7), (8)].

Wir nennen daher die $\sigma_a$ die \gesperrt{Substitutionen des Körpers} $\Omega(\varrho)$.

Wenn $\omega = \varphi(\varrho)$ im Körper $\Omega(\varrho)$ ungeändert bleibt, wenn $\varrho$ durch $\varrho_a$ ersetzt wird, wenn also
%CONT>>
